\section{Related Work}
\label{sec:related}

\paragraph{LLM Routing and Cascading.}
FrugalGPT \citep{chen2023frugalgpt} pioneered cost-aware LLM selection through cascading strategies, achieving 50--90\% cost reduction. Hybrid LLM \citep{ding2024hybridllm} routes queries between small and large models based on predicted difficulty, reducing large-model calls by up to 40\%. RouteLLM \citep{ong2024routellm} trained pairwise classifiers for binary routing using preference data from LMSYS Chatbot Arena. AutoMix \citep{madaan2023automix} used self-verification to decide when to escalate from small to large models. \citet{dekoninck2024cascade} unified routing and cascading into a single framework, showing that iterative model selection with optional skipping outperforms pure routing or pure cascading by up to 14\%. BEST-Route \citep{ding2025bestroute} extended this by jointly selecting both model and number of responses based on query difficulty. RouterBench \citep{hu2024routerbench} and RouterEval \citep{huang2024routereval} provide standardized evaluation frameworks. Most of these systems focus on LLM-to-LLM routing with additive or sequential scoring.

\paragraph{Learned Representations for Routing.}
A growing body of work learns query and model representations for routing decisions. RouteLLM's strongest baseline uses matrix factorization over preference data. RouterDC \citep{chen2024routerdc} trains a router via dual contrastive learning (sample-LLM and sample-sample losses) to match queries with the best-performing LLM from a candidate pool. EmbedLLM \citep{zhuang2024embedllm} learns compact vector embeddings of LLMs using an encoder-decoder approach, enabling performance forecasting across 112+ models. GraphRouter \citep{feng2024graphrouter} constructs heterogeneous task-model graphs and uses edge prediction for routing. ICL-Router \citep{wang2025iclrouter} uses in-context vectors to represent model capabilities, enabling integration of new models without retraining. Prompt-to-Leaderboard \citep{frick2025p2l} trains an LLM to output Bradley-Terry coefficients from prompts, predicting human preferences for prompt-specific model ranking. ZeroRouter \citep{yan2026zerorouter} creates a universal, model-agnostic latent space for query difficulty, enabling zero-shot onboarding of new models using only anchor queries. Our BilinearPhiEngine (Section~\ref{sec:learned_phi}) shares the learned-embedding approach with these systems. The specific contribution is integrating bilinear forms as drop-in replacements for the hand-designed phi functions within the S(M,T) scoring architecture, preserving the framework's multiplicative gating and convergence properties while bypassing the measurement gap.

\paragraph{Multi-Agent and Heterogeneous Routing.}
Several recent systems extend routing beyond LLM-to-LLM selection. MoMA \citep{moma2025routing} integrates LLM-as-a-judge with Mixture-of-Experts for unified routing across both models and agents. DAAO \citep{daao2025} uses a VAE for difficulty estimation and dynamically generates query-specific multi-agent workflows. xRouter \citep{qian2025xrouter} frames routing as a tool-calling problem and trains end-to-end with reinforcement learning. Router-R1 \citep{zhang2025routerr1} teaches an LLM to interleave reasoning and routing actions for multi-round aggregation. S(M,T) also routes across heterogeneous endpoints (LLMs, agents, scripts, tool-use systems), but through a different mechanism: a single scoring equation with endpoint-type-specific gates rather than an LLM-based orchestrator.

\paragraph{Mixture of Experts.}
Sparse Mixture-of-Experts \citep{shazeer2017outrageously} routes tokens to specialized sub-networks within a single model. Switch Transformers \citep{fedus2022switch} simplified gating to top-1 expert selection. These approaches operate at the token level within a model, while S(M,T) operates at the query level across models. The Product-of-Experts framework \citep{hinton2002poe}, which we use for the gate layer, was originally proposed for combining generative models.

\paragraph{Online Learning for Model Selection.}
The contextual bandit formulation for model selection has been explored by \citet{agrawal2013thompson} in the general setting. TI-UCB \citep{tiucb2024} adapts upper confidence bound methods for online model selection with time-increasing reward signals. BaRP \citep{barp2025} frames routing as a multi-objective contextual bandit conditioned on trade-off vectors, enabling performance-cost preference adjustment at inference without retraining. OmniRouter \citep{mei2025omnirouter} models routing as constrained optimization via Lagrangian dual decomposition under budget constraints. Robbins-Monro stochastic approximation \citep{robbins1951stochastic} and convergence under Hajek conditions \citep{hajek1988cooling} are well-established. Soft Actor-Critic \citep{haarnoja2018sac} introduced entropy-regularized learning for continuous control. We combine Robbins-Monro weight updates with SAC-style entropy decay for routing, with Duchi simplex projection \citep{duchi2008simplex} to maintain valid weight distributions.

\paragraph{Benchmark Reliability.}
Concerns about benchmark contamination \citep{deng2023leaderboards}, evaluation methodology \citep{liang2022helm}, and the gap between benchmark performance and real-world utility have been raised repeatedly. Our measurement gap analysis (Section~\ref{sec:measurement}) provides a systematic documentation of how these concerns specifically impact routing, demonstrating 0/14 successful measurement designs across 2.76M records.

\paragraph{What S(M,T) Adds.}
Given the breadth of recent routing work, we clarify the specific contributions of this paper. (1) \emph{A unified scoring equation with formal guarantees}: the three-part multiplicative form (gate products $\times$ weighted compatibility $\times$ Boltzmann cost) with provable convergence ($O(\sqrt{KN\log N})$ regret), parameter identifiability, and adversarial robustness. No prior system provides all three guarantees for the complete routing pipeline. (2) \emph{The measurement gap as an empirical finding}: showing that hand-designed, interpretable scoring functions cannot be grounded in public benchmark data due to metric incomparability, domain transfer breakdown, and dimensional collapse. This finding is independent of any particular routing architecture. (3) \emph{Heterogeneous endpoint support via type-specific gates}: while MoMA and DAAO also route beyond LLMs, S(M,T) handles heterogeneous endpoints through the gate layer within a single equation rather than through an LLM orchestrator, preserving scoring interpretability and convergence properties. (4) \emph{A learned resolution that preserves framework structure}: the BilinearPhiEngine replaces hand-designed phi functions with learned bilinear forms within the same S(M,T) architecture. The bilinear form itself builds on prior work in learned routing representations (RouterDC, EmbedLLM, GraphRouter). The contribution is demonstrating that this approach bypasses the measurement gap while retaining the framework's multiplicative gating and weight-learning machinery.
